\documentclass[11pt,twocolumn]{article}

% --- Packages ---
\usepackage[utf8]{inputenc}
\usepackage[T1]{fontenc}
\usepackage{lmodern}
\usepackage[margin=1in]{geometry}
\usepackage{amsmath,amssymb,amsthm}
\usepackage{algorithm}
\usepackage{algpseudocode}
\usepackage{graphicx}
\usepackage{booktabs}
\usepackage{hyperref}
\usepackage{cleveref}
\usepackage{xcolor}
\usepackage{listings}
\usepackage{caption}
\usepackage{subcaption}
\usepackage{tikz}
\usetikzlibrary{arrows.meta,positioning,shapes.geometric}

% --- Listings style ---
\lstset{
  basicstyle=\ttfamily\small,
  keywordstyle=\bfseries,
  commentstyle=\itshape\color{gray},
  frame=single,
  breaklines=true,
  columns=fullflexible,
}

% --- Theorem environments ---
\newtheorem{theorem}{Theorem}
\newtheorem{lemma}[theorem]{Lemma}
\newtheorem{proposition}[theorem]{Proposition}
\newtheorem{corollary}[theorem]{Corollary}
\theoremstyle{definition}
\newtheorem{definition}{Definition}

% --- Title ---
\title{%
  WISP-0: A Bandwidth-Efficient Protocol for Incremental\\
  Binary Synchronization via Content-Defined Chunking,\\
  Invertible Count Tables, and Fountain Coding%
}
\author{%
  WISP Protocol Working Group\\
  \texttt{wisp-protocol@example.org}
}
\date{March 2026}

\begin{document}
\maketitle

% ===========================================================================
\begin{abstract}
We present WISP-0 (Wire-level Incremental Sync Protocol), a novel
application-layer protocol that achieves bandwidth-efficient incremental
synchronization of opaque binary objects by composing three complementary
techniques: (1)~content-defined chunking (CDC) for edit-local decomposition,
(2)~Invertible Count Tables (ICT) for single-round-trip set reconciliation,
and (3)~systematic fountain coding (SSX) for rateless erasure-resilient
bulk transfer. Control signaling flows over a reliable channel (TCP) while
bulk data flows over an unreliable datagram channel (UDP), avoiding
head-of-line blocking. We define the CSR/0 profile with concrete algorithm
parameters and provide a complete Rust reference implementation
(${\sim}7{,}100$ lines of code) that is \texttt{no\_std}-compatible for
embedded and WebAssembly deployment, with optional X25519-PSK-HKDF
authentication providing per-session forward secrecy. Empirical evaluation
demonstrates that WISP-0 transfers only the bytes that actually changed,
achieves decoding overhead within 5--10\% of the information-theoretic
minimum, and completes synchronization in a fixed number of round trips
regardless of object size.
\end{abstract}

\textbf{Keywords:} incremental synchronization, content-defined chunking,
set reconciliation, invertible bloom lookup tables, fountain codes,
rateless erasure codes, binary delta transfer


% ===========================================================================
\section{Introduction}
\label{sec:intro}

Efficient synchronization of binary objects across distributed systems is a
fundamental problem in systems engineering. Applications range from firmware
distribution on IoT devices, to container image layer synchronization, to
replicated database snapshots and machine learning model deployment.

Existing approaches fall into two broad categories:

\paragraph{Delta Compression (rsync-family).}
The rsync algorithm~\cite{rsync} and its derivatives compute rolling
checksums to identify matching blocks between source and target, then
transfer only non-matching regions. However, rsync requires $O(\log n)$
round trips for checksum negotiation in the general case, and its
fixed-block decomposition causes $O(n)$ block invalidation for a single
byte insertion.

\paragraph{Content-Addressed Transfer (IPFS-family).}
Systems like IPFS~\cite{ipfs} decompose objects into content-addressed
blocks and transfer missing blocks. While this achieves deduplication,
the set reconciliation step typically requires exchanging full block
lists or multiple Bloom filter rounds.

\medskip

WISP-0 improves upon both approaches by combining:

\begin{enumerate}
  \item \textbf{Content-Defined Chunking} (CDC): A rolling-hash chunker
    that places boundaries at content-determined positions, ensuring
    $O(1)$ edit locality---a single-byte edit affects at most a constant
    number of chunks.

  \item \textbf{Invertible Count Tables} (ICT): A compact probabilistic
    sketch that encodes an entire chunk-id set, enabling exact set
    difference recovery in a \emph{single} round trip.

  \item \textbf{Systematic Fountain Coding} (SSX): A rateless erasure
    code over $\text{GF}(2)$ that provides graceful degradation under
    arbitrary packet loss, eliminating per-packet acknowledgment.
\end{enumerate}

The result is a protocol that transfers \emph{exactly} the missing bytes
(plus coding overhead), completes in a fixed number of round trips
(typically 3--4 on C0), and tolerates arbitrary loss patterns on the
bulk data channel without retransmission.

\paragraph{Contributions.}
\begin{itemize}
  \item A complete protocol specification (WISP-0 CSR/0) suitable for
    IETF standardization (\cref{sec:protocol}).
  \item Formal analysis of bandwidth optimality and round-trip complexity
    (\cref{sec:analysis}).
  \item A production-quality \texttt{no\_std} Rust implementation with
    142 tests (unit, integration, and end-to-end) and conformance vectors
    (\cref{sec:impl}).
  \item A defense-in-depth security layer with X25519-PSK-HKDF key
    exchange, HMAC-SHA256-128 frame authentication, replay protection,
    and adaptive timeout management (\cref{sec:security}).
  \item Empirical evaluation on synthetic workloads
    (\cref{sec:evaluation}).
\end{itemize}


% ===========================================================================
\section{Related Work}
\label{sec:related}

\subsection{Delta Synchronization}

The rsync algorithm~\cite{rsync} pioneered rolling-checksum-based delta
transfer. Subsequent work improved upon rsync's fixed-block model:
zsync~\cite{zsync} moved computation to the client side,
casync~\cite{casync} combined CDC with content-addressed storage, and
various cloud storage systems use proprietary CDC variants.

\subsection{Set Reconciliation}

Minsky, Trachtenberg, and Zippel~\cite{mtz} introduced polynomial
interpolation for set reconciliation. Eppstein et al.~\cite{eppstein}
proposed strata estimators with Invertible Bloom Lookup Tables (IBLTs)
for practical single-round reconciliation. Goodrich and
Mitzenmacher~\cite{iblt} formalized IBLTs and analyzed their
performance. Our ICT variant uses fixed $k=3$ hash functions with
SHA-256-derived indices, trading generality for implementation
simplicity and reproducibility.

\subsection{Erasure Coding}

LT codes~\cite{lt-codes} introduced the fountain coding paradigm.
Raptor codes~\cite{raptor} improved upon LT codes with a pre-coding
step for linear-time encoding/decoding. WISP-0's SSX/0 uses a simpler
sparse systematic approach with GF(2) Gaussian elimination, which is
sufficient for the symbol counts typical in file synchronization
($n < 10^5$) and avoids the complexity of Raptor's pre-code design.


% ===========================================================================
\section{Protocol Design}
\label{sec:protocol}

\subsection{Architecture}

WISP-0 employs a dual-channel architecture:

\begin{itemize}
  \item \textbf{C0 (Control)}: TCP, carrying session management,
    recipe exchange, and synchronization frames.
  \item \textbf{D0 (Data)}: UDP, carrying source symbols (SYM) and
    coded repair symbols (MIX).
\end{itemize}

This separation avoids TCP's head-of-line blocking on the bulk data
path while maintaining reliable delivery for protocol correctness.

\subsection{Object Model}

An \emph{Object} is an opaque byte sequence. Each Object version is
decomposed into a sequence of \emph{Chunks} by CDC, described by a
\emph{Recipe}---an ordered list of $(\text{ChunkId}, \text{chunk\_len})$
pairs. The \emph{Recipe Digest} $H_{\text{recipe}} =
\text{SHA-256}(\text{RecipeWire})$ uniquely identifies each version.

\subsection{Protocol Phases}

An update cycle comprises five phases:

\begin{enumerate}
  \item \textbf{Session Establishment}: HELLO exchange + SELECT
    (capability negotiation).
  \item \textbf{Recipe Exchange}: OFFER with RecipeOps delta, ACK.
  \item \textbf{Set Reconciliation}: HAVE\_ICT sketch, ICT peeling.
  \item \textbf{Bulk Transfer}: PLAN, then SYM + MIX on D0.
  \item \textbf{Commit}: COMMIT with digest verification, ACK.
\end{enumerate}

\subsection{Frame Format}

Every frame carries an 80-byte header (see \cref{tab:header}).

\begin{table}[t]
\centering
\caption{WISP-0 Frame Header (80 bytes, little-endian).}
\label{tab:header}
\small
\begin{tabular}{@{}rrlp{3.2cm}@{}}
\toprule
Offset & Size & Field & Description \\
\midrule
0  & 2  & \texttt{magic}        & 0x5057 \\
2  & 1  & \texttt{version}      & Protocol version (0) \\
3  & 1  & \texttt{frame\_type}  & Frame type code \\
4  & 2  & \texttt{flags}        & Reserved (0) \\
6  & 4  & \texttt{payload\_len} & Payload byte count \\
10 & 8  & \texttt{session\_id}  & Session identifier \\
18 & 8  & \texttt{object\_id}   & Object identifier \\
26 & 8  & \texttt{update\_id}   & Update sequence \\
34 & 8  & \texttt{frame\_id}    & Control frame seq\# \\
42 & 4  & \texttt{capsule\_id}  & Profile identifier \\
46 & 1  & \texttt{chan}          & 0=control, 1=data \\
47 & 1  & \texttt{rsv0}         & Reserved (0) \\
48 & 32 & \texttt{base\_digest} & SHA-256 of base recipe \\
\bottomrule
\end{tabular}
\end{table}


% ===========================================================================
\section{Content-Defined Chunking (CSR/0)}
\label{sec:cdc}

\subsection{Rolling Hash}

CSR/0 uses a Rabin-variant polynomial rolling hash over a sliding window
of $W = 48$ bytes:

\begin{equation}
  h_i = h_{i-1} \cdot B + x_i - y_i \cdot B^W \pmod{2^{64}}
  \label{eq:rolling-hash}
\end{equation}

where $B = 257$, $x_i$ is the incoming byte, $y_i$ is the outgoing byte
from the ring buffer, and $B^W \bmod 2^{64}$ is a precomputed constant.

\subsection{Boundary Detection}

A chunk boundary is declared at position $i$ when:

\begin{equation}
  \bigl(\text{chunk\_size} \geq \text{MIN}\bigr) \wedge
  \bigl(h_i \mathbin{\&} \text{MASK} = 0\bigr)
  \label{eq:boundary}
\end{equation}

or when $\text{chunk\_size} = \text{MAX}$ (forced cut).

\begin{table}[t]
\centering
\caption{CSR/0 Chunker Parameters.}
\label{tab:cdc-params}
\begin{tabular}{@{}lrl@{}}
\toprule
Parameter & Value & Description \\
\midrule
$B$    & 257      & Hash base \\
$W$    & 48       & Window size (bytes) \\
$K$    & 16       & Mask bit width \\
MIN    & 8\,192   & Minimum chunk (8\,KiB) \\
MAX    & 262\,144 & Maximum chunk (256\,KiB) \\
MASK   & 0xFFFF   & $(1 \ll K) - 1$ \\
\bottomrule
\end{tabular}
\end{table}

\begin{proposition}[Edit Locality]
\label{prop:locality}
A modification of $\delta$ consecutive bytes in the input affects at most
$\lceil \delta / \text{MIN} \rceil + 2$ chunks. Chunks whose content is
entirely outside the modified region retain identical ChunkIds.
\end{proposition}

\begin{proof}
The rolling hash window extends $W = 48$ bytes. A modification at
position $p$ can only affect boundary decisions within the range
$[p - W, p + \delta + \text{MAX}]$. Since MIN $= 8192 \gg W$, the
perturbation is confined to at most
$\lceil (\delta + 2 \cdot \text{MAX}) / \text{MIN} \rceil$ chunk
boundary decisions. In practice, the two boundary chunks adjacent to the
modification may shift, contributing the $+2$ term, while interior chunks
are completely determined by the modified content.
\end{proof}


% ===========================================================================
\section{Set Reconciliation (ICT/0)}
\label{sec:ict}

\subsection{Data Structure}

An Invertible Count Table with $m = 2^{m_{\log_2}}$ cells, each
containing:

\begin{equation}
  \text{cell} = (\text{count} \in \mathbb{Z},\
                 \text{key\_xor} \in \{0,1\}^{128},\
                 \text{hash\_xor} \in \{0,1\}^{64})
\end{equation}

\subsection{Hash Functions}

ICT/0 uses $k = 3$ hash functions derived from SHA-256 with domain
separation:

\begin{align}
  \text{hash64}(\text{key}) &= \text{LE}_{64}\bigl(
    \text{SHA-256}(0\text{x}49 \| \text{key})[0{:}8]\bigr) \\
  h_j(\text{key}, m) &= \text{LE}_{64}\bigl(
    \text{SHA-256}((0\text{xA0}{+}j) \| \text{key})[0{:}8]\bigr)
    \bmod m
\end{align}

for $j \in \{0, 1, 2\}$.

\subsection{Table Sizing}

Given an estimate $d$ of the symmetric difference:

\begin{equation}
  m = \max\bigl(\text{next\_pow2}\bigl(\lceil 2.5d + 32 \rceil\bigr),
  256\bigr)
  \label{eq:sizing}
\end{equation}

This provides $\approx 2.5\times$ overhead, ensuring high peeling success
probability.

\begin{theorem}[Peeling Success]
\label{thm:peel}
For $k = 3$ hash functions mapping uniformly into $m$ cells, the peeling
algorithm recovers all $d$ elements from the difference table with
probability $\geq 1 - m^{-\Omega(1)}$ when $m \geq 2.5d + 32$.
\end{theorem}

\begin{proof}[Proof sketch]
The analysis follows from the random hypergraph peelability threshold for
$k$-uniform hypergraphs~\cite{iblt}. For $k = 3$, the threshold ratio
$m/d$ approaches $\approx 1.222$ as $d \to \infty$. Our sizing formula
provides ratio $\geq 2.5$, well above threshold, yielding exponentially
small failure probability. The additive constant 32 ensures adequate
table size for small $d$.
\end{proof}

\subsection{Operations}

\begin{algorithm}[t]
\caption{ICT Insert and Peel}
\label{alg:ict}
\begin{algorithmic}[1]
\Procedure{Insert}{$\text{table}, \text{key}$}
  \State $h \gets \text{hash64}(\text{key})$
  \For{$j = 0, 1, 2$}
    \State $i \gets h_j(\text{key}, m)$
    \State $\text{table}[i].\text{count} \mathrel{+}= 1$
    \State $\text{table}[i].\text{key\_xor} \mathrel{\oplus}= \text{key}$
    \State $\text{table}[i].\text{hash\_xor} \mathrel{\oplus}= h$
  \EndFor
\EndProcedure
\Statex
\Procedure{Peel}{$\text{table}$}
  \State $\text{stack} \gets \{i : \text{is\_pure}(\text{table}[i])\}$
  \State $\text{missing} \gets \emptyset$
  \While{$\text{stack} \neq \emptyset$}
    \State $i \gets \text{stack.pop}()$
    \If{$\lnot\text{is\_pure}(\text{table}[i])$} \textbf{continue} \EndIf
    \State $\text{key} \gets \text{table}[i].\text{key\_xor}$
    \State $\text{sign} \gets \text{table}[i].\text{count}$
    \If{$\text{sign} = +1$} $\text{missing} \gets \text{missing}
      \cup \{\text{key}\}$ \EndIf
    \State Remove $\text{key}$ from all 3 cell positions
    \State Push newly pure cells onto stack
  \EndWhile
  \State \Return missing
\EndProcedure
\end{algorithmic}
\end{algorithm}


% ===========================================================================
\section{Fountain Coding (SSX/0)}
\label{sec:ssx}

\subsection{Source Symbol Construction}

Missing chunks are serialized into $n$ source symbols of size
$L = 1024$ bytes. Each chunk is split into $\lceil \text{len}/L \rceil$
symbols, zero-padded to $L$ bytes.

\subsection{Degree Distribution}

The SSX/0 degree distribution is determined by a single PRNG byte $b$:

\begin{table}[t]
\centering
\caption{SSX/0 Degree Distribution.}
\label{tab:degree}
\begin{tabular}{@{}ccr@{}}
\toprule
Byte Range & Degree $d$ & Probability \\
\midrule
$[0, 128)$   & 3  & 50.0\% \\
$[128, 192)$ & 5  & 25.0\% \\
$[192, 224)$ & 9  & 12.5\% \\
$[224, 256)$ & 17 & 12.5\% \\
\bottomrule
\end{tabular}
\end{table}

The expected degree is:

\begin{equation}
  \mathbb{E}[d] = 0.5 \times 3 + 0.25 \times 5 + 0.125 \times 9 +
  0.125 \times 17 = 5.0
\end{equation}

\subsection{PRNG and Seed Derivation}

The PRNG is a SHA-256 hash chain initialized with:

\begin{equation}
  \text{seed}_t = \text{SHA-256}\bigl(\text{``SSX0''} \|
  \text{sid}_{64} \| \text{oid}_{64} \| \text{uid}_{64} \|
  t_{32}\bigr)
\end{equation}

where all integers are encoded as little-endian bytes.

\subsection{Encoding}

A coded symbol for fountain index $t$ is:

\begin{equation}
  c_t = \bigoplus_{i \in \mathcal{I}_t} s_i
\end{equation}

where $\mathcal{I}_t$ is the set of $d_t$ indices drawn from the PRNG
and $\oplus$ denotes byte-wise XOR.

\subsection{Decoding}

The decoder maintains a sparse system of linear equations over
$\text{GF}(2)$. Each received symbol (SYM or MIX) adds one equation.
Known symbols are substituted into pending equations, and
single-variable equations are immediately resolved (peeling).

\begin{theorem}[Decoding Overhead]
\label{thm:overhead}
With the SSX/0 degree distribution (expected degree 5.0), a receiver
that collects $n(1 + \varepsilon)$ symbols (systematic + coded)
recovers all $n$ source symbols with probability approaching 1 as
$n \to \infty$, for any $\varepsilon > 0$.
\end{theorem}

\begin{proof}[Proof sketch]
The analysis parallels the density evolution argument for LT
codes~\cite{lt-codes}. The bipartite graph between $n$ source symbols
and $n(1+\varepsilon)$ coded symbols has expected degree 5.0 on the
coded side. By the theory of random bipartite graphs, the iterative
peeling process (equivalent to belief propagation on the binary erasure
channel) succeeds with high probability when the degree distribution
lies above the BP threshold for the erasure rate
$\delta = \varepsilon / (1 + \varepsilon)$. The Gaussian elimination
fallback handles the $O(\sqrt{n})$ residual symbols that resist peeling.
\end{proof}


% ===========================================================================
\section{Analysis}
\label{sec:analysis}

\subsection{Bandwidth Optimality}

\begin{definition}[Bandwidth overhead]
For an update where $\Delta$ bytes of chunk content differ between base
and target, the \emph{bandwidth overhead} is
$\rho = (\text{bytes transferred}) / \Delta$.
\end{definition}

\begin{theorem}[Asymptotic Optimality]
\label{thm:bandwidth}
WISP-0 with CSR/0 achieves bandwidth overhead
$\rho = 1 + O(\varepsilon)$ as $\Delta \to \infty$, where $\varepsilon$
is the fountain coding overhead.
\end{theorem}

\begin{proof}
The protocol transfers: (1)~RecipeOps of size
$O(|\text{changed chunks}| \cdot 18)$ bytes on C0, (2)~ICT sketch of
size $O(d \cdot 28 \cdot 2.5)$ bytes on C0, (3)~$n(1+\varepsilon)$
fountain symbols of $L = 1024$ bytes on D0, where
$n = \lceil \Delta / L \rceil$. The dominant term is
$(3) = \Delta(1 + \varepsilon) + O(L)$, giving
$\rho = 1 + \varepsilon + O(1/n)$.
\end{proof}

\subsection{Round-Trip Complexity}

\begin{proposition}[Fixed Round Trips]
A WISP-0 update cycle on C0 requires exactly 4 round trips in the
common case (no NEEDMORE):
\begin{enumerate}
  \item HELLO exchange (session setup, amortized)
  \item OFFER $\to$ ACK
  \item HAVE\_ICT $\to$ PLAN $\to$ ACK
  \item COMMIT $\to$ ACK
\end{enumerate}
\end{proposition}

This is independent of object size, in contrast to rsync's
$O(\log n)$ round trips.

\subsection{ICT Communication Complexity}

The ICT sketch has size $28 \cdot m = 28 \cdot 2^{m_{\log_2}}$ bytes.
For $d$ missing chunks:

\begin{equation}
  |\text{ICT}| = 28 \cdot 2^{\lceil \log_2(2.5d + 32) \rceil}
  \leq 28 \cdot 2 \cdot (2.5d + 32)
  = O(d)
\end{equation}

This is linear in the size of the difference, compared to $O(n)$ for
exchanging full chunk lists.


% ===========================================================================
\section{Implementation}
\label{sec:impl}

\subsection{Architecture}

The reference implementation comprises four Rust crates:

\begin{table}[t]
\centering
\caption{Crate structure of the reference implementation.}
\label{tab:crates}
\begin{tabular}{@{}llr@{}}
\toprule
Crate & Purpose & LOC \\
\midrule
\texttt{wisp-core}    & Protocol logic (\texttt{no\_std}) & $\sim$4\,460 \\
\texttt{wisp-net}     & Async networking (Tokio)           & $\sim$1\,950 \\
\texttt{wisp-cli}     & CLI binary (send/recv)             & $\sim$190 \\
\texttt{wisp-vectors} & Conformance vectors                & $\sim$510 \\
\bottomrule
\end{tabular}
\end{table}

\subsection{\texttt{no\_std} Compatibility}

\texttt{wisp-core} is fully \texttt{no\_std}-compatible with the
\texttt{alloc} crate, enabling deployment on:
\begin{itemize}
  \item Embedded systems (ARM Cortex-M, RISC-V)
  \item WebAssembly (\texttt{wasm32-unknown-unknown})
  \item Kernel-mode drivers
  \item Resource-constrained IoT devices
\end{itemize}

The \texttt{std} feature is on by default and adds
\texttt{thiserror}-derived error types and \texttt{std::io::Error}
integration. Opting out via \texttt{default-features = false} removes
all standard library dependencies.

\subsection{Test Coverage}

The implementation includes 143 tests:
\begin{itemize}
  \item 119 unit tests across 12 modules
  \item 7 integration tests (incremental sync, loopback, conformance)
  \item 10 end-to-end async networking tests (including incremental delta transfer)
  \item 7 RTT/pacer/session unit tests
\end{itemize}

All tests pass in both \texttt{std} and \texttt{no\_std} configurations
with zero compiler warnings.

\subsection{Conformance Vectors}

Four sets of conformance vectors (JSON format) cover all protocol
components: CDC chunking, Recipe serialization, ICT operations, and
SSX fountain coding. These enable interoperability testing for
independent implementations.


% ===========================================================================
\section{Security}
\label{sec:security}

WISP-0 provides an optional defense-in-depth security layer that can be
negotiated during session establishment. When enabled, the protocol
delivers per-session forward secrecy, frame-level integrity, and replay
protection without relying on an external TLS or DTLS wrapper.

\subsection{Authentication Model}

Key agreement follows an X25519-PSK-HKDF scheme~\cite{rfc7748,rfc5869}.
During the HELLO exchange, both peers include ephemeral X25519 public keys.
Each peer computes the raw shared secret
$\mathit{ss} = \text{X25519}(\mathit{eph\_sk}, \mathit{eph\_pub\_remote})$
and derives the session key as:

\begin{equation}
  k_{\text{session}} = \text{HKDF-SHA256}\!\bigl(
    \mathit{salt} = \mathit{session\_id},\;
    \mathit{ikm} = \mathit{ss} \| \mathit{PSK},\;
    \mathit{info} = \text{``wisp0-session''}
  \bigr)
  \label{eq:session-key}
\end{equation}

The optional pre-shared key (PSK) binds the session to a mutually known
secret, preventing man-in-the-middle attacks when the PSK is distributed
out of band. Ephemeral key pairs ensure forward secrecy: compromise of
the PSK alone does not expose past session keys.

\subsection{Frame Authentication}

Authenticated frames carry an HMAC-SHA256-128 tag (16~bytes) appended
after the payload. The tag is computed over
$\texttt{header} \| \texttt{payload}$ using $k_{\text{session}}$.
Bit~0 of the \texttt{flags} header field signals that a frame is
authenticated; receivers reject tagged frames whose HMAC does not verify
and silently discard untagged frames when authentication is required by
session policy.

\subsection{Replay Protection}

D0 frames are protected by a 128-bit sliding window indexed by the
per-frame \texttt{frame\_id}. Frames whose identifier falls outside or
has already been recorded in the window are discarded. C0 frames enforce
monotonically increasing \texttt{frame\_id} values; any non-monotonic
control frame is treated as a protocol violation and terminates the
session.

\subsection{Operational Hardening}

Three mechanisms improve robustness under adverse network conditions:
\begin{itemize}
  \item \textbf{EWMA RTT estimator}: Smoothed round-trip time and RTT
    variance are maintained per RFC~6298, driving retransmission timeouts
    and the D0 send pacer.
  \item \textbf{D0 send pacer}: Limits the outbound symbol rate to avoid
    congesting the datagram path, adapting to the estimated RTT.
  \item \textbf{NEEDMORE stall recovery}: When the fountain decoder
    cannot make progress, the receiver issues a NEEDMORE frame on C0,
    prompting the sender to transmit additional coded symbols.
\end{itemize}

\subsection{Session Limits}

Each session enforces configurable bounds on total frame count and
elapsed duration. Exceeding either limit causes graceful session
termination, bounding resource consumption under prolonged or
misbehaving connections.


% ===========================================================================
\section{Evaluation}
\label{sec:evaluation}

\subsection{Experimental Setup}

We evaluate WISP-0 on synthetic workloads using the reference
implementation. Test scenarios:

\begin{enumerate}
  \item \textbf{Small edit}: 200\,KB file, 40\,KB middle region modified.
  \item \textbf{Append}: 100\,KB file extended by 50\,KB.
  \item \textbf{Identity}: Same file on both sides (zero transfer).
\end{enumerate}

\subsection{Bandwidth Efficiency}

\begin{table}[t]
\centering
\caption{Transfer efficiency on synthetic workloads.}
\label{tab:efficiency}
\begin{tabular}{@{}lrrrr@{}}
\toprule
Scenario & $|\text{File}|$ & $\Delta$ & Transferred & $\rho$ \\
\midrule
Small edit & 200\,KB & 40\,KB  & 42\,KB & 1.05 \\
Append     & 150\,KB & 50\,KB  & 53\,KB & 1.06 \\
Identity   & 80\,KB  & 0       & 0.3\,KB & --- \\
\bottomrule
\end{tabular}
\end{table}

The identity case transfers only protocol overhead (OFFER with empty
RecipeOps + ICT sketch with zero difference).

\subsection{Fountain Coding Performance}

The SSX/0 decoder consistently achieves complete recovery with
5--10\% overhead beyond the systematic symbols:

\begin{itemize}
  \item With 25\% systematic symbols (SYM) + coded symbols (MIX),
    recovery completes within $\sim 2n$ total symbols.
  \item Pure MIX recovery (no SYM) requires $\sim 1.1n$ symbols.
  \item The Gaussian elimination fallback typically resolves
    $< 5\%$ of symbols that resist iterative peeling.
\end{itemize}

\subsection{ICT Reconciliation}

The ICT peel succeeds deterministically for all test cases with the
$2.5\times$ sizing formula. For $d = 200$ missing chunks, the ICT
sketch is $\approx 14\,\text{KB}$ (512 cells $\times$ 28 bytes),
transmitted in a single HAVE\_ICT frame.


% ===========================================================================
\section{Discussion}
\label{sec:discussion}

\subsection{Comparison with rsync}

WISP-0 improves upon rsync in three key dimensions:

\begin{enumerate}
  \item \textbf{Round trips}: WISP-0 requires $O(1)$ round trips
    (typically 4) versus rsync's $O(\log n)$.
  \item \textbf{Edit locality}: CDC ensures $O(1)$ affected chunks
    per edit versus rsync's $O(n)$ with fixed blocks.
  \item \textbf{Loss resilience}: Fountain coding eliminates
    per-packet ACK, enabling high throughput over lossy links.
\end{enumerate}

\subsection{Limitations}

\begin{itemize}
  \item The ICT sketch size is proportional to the difference size $d$,
    not the object size $n$. When $d \approx n$ (completely different
    files), WISP-0 degenerates to a full transfer with ICT overhead.
  \item The SSX/0 degree distribution is not information-theoretically
    optimal (unlike Raptor codes). For very large $n$, a pre-coded
    design would reduce overhead.
  \item The current implementation does not support streaming chunking
    (the entire object must be available for CDC).
\end{itemize}

\subsection{Future Work}

\begin{itemize}
  \item \textbf{TLS/DTLS wrapping for confidentiality}: WISP-0's native
    HMAC-based authentication (\cref{sec:security}) provides integrity
    and authenticity but does not encrypt payload content. An optional
    TLS (C0) or DTLS (D0) wrapper would add confidentiality where
    required.
  \item \textbf{Formal verification}: Machine-checked proofs of
    protocol correctness and the security layer using Verus or Prusti
    remain desirable for high-assurance deployments.
  \item \textbf{Streaming CDC}: Support incremental chunking as bytes
    arrive, enabling pipeline parallelism.
  \item \textbf{Raptor pre-coding}: Add an LDPC pre-code for
    linear-time decoding at scale.
  \item \textbf{Multi-object sessions}: Support concurrent
    synchronization of multiple objects within a single session.
\end{itemize}


% ===========================================================================
\section{Conclusion}
\label{sec:conclusion}

We have presented WISP-0, a complete protocol for bandwidth-efficient
incremental binary synchronization. By composing content-defined
chunking, invertible count tables, and fountain coding, WISP-0 achieves
near-optimal bandwidth utilization with a fixed number of round trips
and graceful degradation under packet loss. The CSR/0 profile provides
concrete, reproducible algorithm parameters, and the reference
implementation demonstrates practical deployability across platforms
from embedded devices to cloud infrastructure. An optional X25519-PSK-HKDF
security layer provides per-session forward secrecy and frame-level
integrity without introducing external dependencies.

The protocol's modular design---with pluggable chunker profiles, ICT
parameters, and fountain coding schemes---provides a foundation for
future evolution while maintaining backward compatibility through the
capsule identifier mechanism.


% ===========================================================================
\bibliographystyle{plain}
\begin{thebibliography}{99}

\bibitem{rsync}
A.~Tridgell and P.~Mackerras,
``The rsync algorithm,''
Technical Report TR-CS-96-05, Australian National University, 1996.

\bibitem{ipfs}
J.~Benet,
``IPFS -- Content Addressed, Versioned, P2P File System,''
\textit{arXiv:1407.3561}, 2014.

\bibitem{iblt}
M.~T.~Goodrich and M.~Mitzenmacher,
``Invertible Bloom Lookup Tables,''
in \textit{Proc.\ 49th Annual Allerton Conference}, 2011.

\bibitem{lt-codes}
M.~Luby,
``LT Codes,''
in \textit{Proc.\ 43rd Annual IEEE Symposium on Foundations of Computer
Science (FOCS)}, 2002, pp.~271--280.

\bibitem{raptor}
A.~Shokrollahi,
``Raptor Codes,''
\textit{IEEE Transactions on Information Theory}, vol.~52, no.~6,
pp.~2551--2567, June 2006.

\bibitem{cdc}
K.~Eshghi and H.~K.~Tang,
``A Framework for Analyzing and Improving Content-Based Chunking
Algorithms,''
HPL-2005-30R1, Hewlett-Packard Labs, 2005.

\bibitem{rabin}
M.~O.~Rabin,
``Fingerprinting by Random Polynomials,''
Center for Research in Computing Technology, Harvard University,
TR-15-81, 1981.

\bibitem{mtz}
Y.~Minsky, A.~Trachtenberg, and R.~Zippel,
``Set Reconciliation with Nearly Optimal Communication Complexity,''
\textit{IEEE Transactions on Information Theory}, vol.~49, no.~9,
pp.~2213--2218, September 2003.

\bibitem{eppstein}
D.~Eppstein, M.~T.~Goodrich, F.~Uyeda, and G.~Varghese,
``What's the Difference? Efficient Set Reconciliation without Prior
Context,''
in \textit{Proc.\ ACM SIGCOMM}, 2011, pp.~218--229.

\bibitem{zsync}
C.~Blake,
``zsync -- optimized rsync over HTTP,''
\url{http://zsync.moria.org.uk/}, 2005.

\bibitem{casync}
L.~Poettering,
``casync -- Content-Addressable Data Synchronization Tool,''
\url{https://github.com/systemd/casync}, 2017.

\bibitem{rfc5869}
H.~Krawczyk and P.~Eronen,
``HMAC-based Extract-and-Expand Key Derivation Function (HKDF),''
RFC~5869, Internet Engineering Task Force, May 2010.

\bibitem{rfc7748}
A.~Langley, M.~Hamburg, and S.~Turner,
``Elliptic Curves for Security,''
RFC~7748, Internet Engineering Task Force, January 2016.

\end{thebibliography}


% ===========================================================================
\appendix

\section{Frame Type Registry}
\label{app:frames}

\begin{table}[h]
\centering
\caption{Complete frame type registry for WISP-0.}
\begin{tabular}{@{}clcc@{}}
\toprule
Code & Name & Chan & Payload Size \\
\midrule
0x01 & HELLO     & 0 & Variable ($\geq$14) \\
0x02 & SELECT    & 0 & 8 \\
0x03 & OFFER     & 0 & Variable ($\geq$32) \\
0x04 & HAVE\_ICT & 0 & Variable \\
0x05 & ICT\_FAIL & 0 & 4 \\
0x06 & PLAN      & 0 & 40 \\
0x07 & NEEDMORE  & 0 & 4 \\
0x08 & COMMIT    & 0 & 32 \\
0x09 & ACK       & 0 & 8 \\
0x0A & FORK      & 0 & 68 \\
0x20 & SYM       & 1 & 1028 \\
0x21 & MIX       & 1 & 1028 \\
\bottomrule
\end{tabular}
\end{table}


\section{RecipeOps/0 Opcode Table}
\label{app:opcodes}

\begin{table}[h]
\centering
\caption{RecipeOps/0 opcodes.}
\begin{tabular}{@{}cll@{}}
\toprule
Code & Name & Parameters \\
\midrule
0x00 & END  & (none) \\
0x01 & COPY & \texttt{base\_index}: ULEB128, \texttt{count}: ULEB128 \\
0x02 & LIT  & \texttt{count}: ULEB128, then \texttt{count} entries \\
\bottomrule
\end{tabular}
\end{table}


\section{Notation Summary}

\begin{table}[h]
\centering
\caption{Symbol notation used throughout this paper.}
\begin{tabular}{@{}ll@{}}
\toprule
Symbol & Meaning \\
\midrule
$n$ & Number of source symbols \\
$L$ & Symbol size in bytes (1024) \\
$d$ & Estimated symmetric difference size \\
$m$ & ICT table size ($2^{m_{\log_2}}$) \\
$k$ & Number of ICT hash functions (3) \\
$B$ & Rolling hash base (257) \\
$W$ & Rolling hash window size (48) \\
$H_{128}$ & Truncated SHA-256 (16 bytes) \\
$H_{256}$ & Full SHA-256 (32 bytes) \\
$\oplus$ & Byte-wise XOR \\
$\rho$ & Bandwidth overhead ratio \\
$\varepsilon$ & Fountain coding overhead fraction \\
\bottomrule
\end{tabular}
\end{table}


\end{document}
